\documentclass{article}
\usepackage[utf8]{inputenc}
\usepackage{fontspec}
\setmainfont[
    Path = ./,
    Ligatures = TeX,
    BoldFont = SABONLTSTD-BOLD.OTF,
    ItalicFont = SABONLTSTD-ITALIC.OTF,
]{SABONLTSTD-ROMAN.OTF}
\usepackage{book-of-common-prayer}
\geometry{
  paperheight=8.5in,
  paperwidth=5.5in,
  heightrounded,
  margin=1in
}

\title{The \textsc{Daily Offices} according to the \textsc{Anglican Breviary}, with the \textsc{Psalter} of King David}
\date{}
\author{}


\begin{document}

\maketitle
\newpage

\tableofcontents

\newpage

\section{General Notes}

\instruct{In the following I put personal notes and not the full rubrics as they are printed in the \textsc{Anglican Breviary}. While the rubrics are good and helpful, they can sometimes be very cumbersome!}

\newpage

\section{Prime}

\subsection*{Prayers before the Office}

\instruct{If \textsc{Prime} is said in sequence with other Offices, then this can be ommitted and said before the first.}

\subsubsection*{Aperi, Domine.}

Open thou, O \textsc{Lord}, my mouth to bless thy holy Name; cleanse also my heart from all vain, evil, and wandering thoughts; enlighten my understanding; enkindle my affections; that I may say this Office worthily with attention and devotion, and so be meet to be heard in the presence of thy divine Majesty. Through Christ our Lord. Amen.

\subsubsection*{The Prayer of St. Gertrude}

O \textsc{Lord}, in union with that divine intention wherewith thou thyself on earth didst render thy praises to God, I desire to offer this my Office of prayer unto thee.

\subsection*{Triple Prayer}

\instruct{The \textsc{Triple Prayer} (The \textsc{Our Father}, \textsc{Hail Mary}, and \textsc{Apostle's Creed}) is always said secretly before the opening versicles of \textsc{Prime}. Even in the case where
the offices are said in sequence, one says this invariably. Typically, these are said standing as one is able. One may make a profound
reverence until the \textsc{Opening Versicles} are read.}

\begin{prayer}

  Our Father, \textit{who art} in heaven,
    hallowed be thy Name.
    Thy kingdom come,
    Thy will be done,
      on earth as it is in heaven.
  Give us this day our daily bread.
  And forgive us our trespasses,
    as we forgive \textit{those who} trespass against us
  And lead us not into temptation,
    but deliver us from evil. Amen.
    
\end{prayer}

\instruct{The \textsc{American Prayerbook} uses \textbf{who art} and \textbf{those who}.}

\instruct{The \textsc{Hail Mary} is composed of the Angelic Salutation to the Blessed Virgin Mary that she would conceive of the Holy Ghost, and where she became the first Christian disciple in \bibleref{Luke 1:28}. It is also composed of her cousin St. Elizabeth's greeting in \bibleref{Luke 1:42}. The final part is a petition for her prayer.}



\begin{prayer}

  Hail Mary, full of grace, the \textsc{Lord} is with thee:
  Blessed art thou among women,
    and blessed is the fruit of thy womb, Jesus.
  Holy Mary, Mother of God, 
    pray for us sinners now, 
    and at the hour of our death. Amen.
    
\end{prayer}

\instruct{Next, is the \textsc{Apostle's Creed}, or \textsc{Credo}.}

\begin{prayer}

  I believe in God, the Father almighty,
    maker of heaven and earth;
  And in Jesus Christ his only Son our Lord;
    who was conceived by the Holy Ghost,
    born of the Virgin Mary,
    suffered under Pontius Pilate,
    was crucified, dead, and buried.
    He descended into hell.
  The third day he rose again from the dead.
    He ascended into heaven,
    and sitteth on the right hand of God the Father almighty.
    From thence he shall come to judge the quick and the dead.
  I believe in the Holy Ghost,
    the holy catholic Church,
    the communion of saints,
    the forgiveness of sins,
    the resurrection of the body,
    and \cross the life everlasting. Amen.
    
\end{prayer}

\newpage

\section{Compline}

\subsection*{Prayers before the Office}

\instruct{If \textsc{Compline} is said in sequence with other Offices, then this can be ommitted and said before the first.}

\subsubsection*{Aperi, Domine.}

Open thou, O \textsc{Lord}, my mouth to bless thy holy Name; cleanse also my heart from all vain, evil, and wandering thoughts; enlighten my understanding; enkindle my affections; that I may say this Office worthily with attention and devotion, and so be meet to be heard in the presence of thy divine Majesty. Through Christ our Lord. Amen.

\subsubsection*{The Prayer of St. Gertrude}

O \textsc{Lord}, in union with that divine intention wherewith thou thyself on earth didst render thy praises to God, I desire to offer this my Office of prayer unto thee.

\subsection*{The Opening Prayers}

\begin{responses}
    \V{Pray, Lord, give me thy blessing.}
    \R{May the Lord almighty grant us a night of quietude, and perfection at the end of all our days.}\
    \textit{Amen.}
\end{responses}

\instruct{Then is said the \textsc{Brief Lesson}, or some other suitable reading of \textsc{Scripture}.}

\bibleverse{Be sober, be vigilant; because your adversary the devil, as a roaring lion, walketh about, seeking whom he may devour: Whom resist stedfast in the faith, knowing that the same afflictions are accomplished in your brethren that are in the world. }{1 Peter 5:8-9}


\instruct{Then is said the \textsc{Lord's Prayer}, secretly.}
\\
%
\begin{prayer}
  Our Father, \textit{who art} in heaven,
    hallowed be thy Name.
    Thy kingdom come,
    Thy will be done,
      on earth as it is in heaven.
  Give us this day our daily bread.
  And forgive us our trespasses,
    as we forgive \textit{those who} trespass against us
  And lead us not into temptation,
    but deliver us from evil. Amen.
    
\end{prayer}

\instruct{The \textsc{American Prayerbook} uses \textbf{who art} and \textbf{those who}.}

\instruct{Then is said the \textsc{Confiteor}, or the \textsc{Confession of Sin} (pp. 62-63, 1979 BCP), audibly.}

\begin{prayer}

  I confess to Almighty God,
    to Blessed Mary Ever-Virgin,
    to Blessed Michael the Archangel,
    to Blessed John the Baptist,
    to the holy Apostles Peter and Paul,
    and to all the Saints,
  that I have sinned sinned exceedingly
    in thought, word, and deed,
  \textit{through my fault, through my own fault,
    through my own most grievous fault}.
  Therefore I beg Blessed Mary Ever-Virgin,
    Blessed Michael the Archangel,
    Blessed John the Baptist,
    the Holy Apostles Peter and Paul,
    and all the Saints,
    to pray to the \textsc{Lord} our God for me.
    
\end{prayer}

\instruct{At \textbf{through my fault, etc.}, strike breast thrice. Then is said.}

Almighty and merciful \textsc{Lord} grant us pardon, \cross absolution, and remission of our sins. Amen.

\begin{responses}

    \V{Turn us then O God our Savior.}
    \R{\cross And let thine anger cease from us.}
    
\end{responses}

\subsection*{The Opening Versicles and Psalms}

\begin{responses}

    \V{O God, make speed to save us.}
    \R{\cross O Lord, make haste to help us.}
    \V{Glory be to the Father, and to the Son, and to the Holy Ghost.}
    \R{As it was in the beginning, is now, and ever shall be, world without end. Amen.}
    
\end{responses}

\instruct{Then is said the appointed \textsc{Psalms} of the Day.}

\subsection*{Hymn}

\instruct{Then is sung the hymn, \textsc{Te lucis ante terminum}, or another suitable one.}

\instruct{For the notation, see \textsc{Supplement}.}

\begin{prayer}

  To thee before the close of day,
  Creator of the world, we pray,
  That with thy wonted favor, thou,
  Wouldst be our Guard and Keeper now.

  From all ill dreams defend our eyes,
  From nightly fears and fantasies:
  Tread under foot our ghostly foe,
  That no pollution we may know.

  O Father, that we ask be done
  Through Jesus Christ, thine only Son,
  Who with the Holy Ghost and thee,
  Shall live and reign eternally. Amen.
    
\end{prayer}

\subsection*{The Little Chapter and Brief Response}

\bibleverse{Thou, O \textsc{Lord}, art in the midst of us, and we are called by thy Name; leave us not, O \textsc{Lord}, our God.}{Jeremiah 14:9}

\textbf{Thanks be to God.}

\begin{responses}
    \V{Into thy hands, O \textsc{Lord}, I commend my spirit.}
    \R{Into thy hands, O \textsc{Lord}, I commend my spirit.}
    \V{For thou hast redeemed us, O \textsc{Lord}, thou God of truth.}
    \R{I commend my spirit.}

    \V{Glory be to the Father, and to the Son, and to the Holy Ghost.}
    \R{Into thy hands, O \textsc{Lord}, I commend my spirit.}
    \V{Keep us, O \textsc{Lord}, as an apple of an eye.}
    \R{Hide us under the shadow of thy wings.}
\end{responses}

\subsection*{The \textsc{Nunc Dimittis} \bibleref{Luke 2:29-32}}

\instruct{Antiphon}

Save us, \textsc{Lord}.

\instruct{Then the prayer is said or sung.}

\begin{prayer}

  Lord, now lettest thou thy servant depart in peace,
    according to thy word:
  For mine eyes have seen thy salvation,
    which thou hast prepared before the face of all people.
  To be a light to lighten the Gentiles,
    and to be the glory of thy people Israel.
\end{prayer}
\blankline
\begin{responses}

    \V{Glory be to the Father, and to the Son, and to the Holy Ghost.}
    \R{As it was in the beginning, is now, and ever shall be, world without end. Amen.}
    
\end{responses}

\subsubsection*{Antiphon}

\textsc{Save} us, \textsc{Lord}, when awake we keep vigil, and guard us when weary we slumber,
that all our watching may be with Christ, and that in him we may rest in peace.

\subsection*{Collect}

\instruct{When a person of Deacon's Orders is present, they say \textbf{The Lord be with you, etc.}. When multiple are gathered, one can substitute \textbf{my} for \textbf{our}.}

\begin{responses}

    \V{Lord, hear my prayer.}
    \R{And let my cry come unto thee.}
    
\end{responses}

Visit, we beseech thee, O Lord, this habitation, and drive far from it all snares of the enemy: let thy holy Angels dwell herein, to preserve us in peace; and let thy blessing be always upon us. Through Jesus Christ, thy son our Lord. Who liveth and reigneth with thee, in the unity of the Holy Ghost, ever one God, world without end.

\begin{responses}
    \R{Amen.}
\end{responses}

\instruct{Likewise, if they have Deacon's orders, they say \textbf{The Lord be with, etc.}}

\begin{responses}
    \V{Lord, hear my prayer.}
    \R{And let my cry come unto thee.}
    \V{Let us bless the \textsc{Lord}.}
    \R{Thanks be to God.}
    \V{May the Lord Almighty and merciful, the Father, \cross the Son, and the Holy Ghost, vouchsafe to bless us and keep us.}
    \R{Amen.}
\end{responses}

\subsection*{Marian Antiphon}

\instruct{The Church has always since Apostolic times asked the Blessed Virgin Mary, Mother of God, to pray for them. In keeping with this tradition, the Office of Compline includes a seasonal hymn to her, petitioning her for her prayers. Thus is said the \textsc{Marian Antiphon} of the season. Below is \textsc{Salve regina}.}

All hail, O holy Queen, Mother exceeding merciful; Life's spring, sweet comfort, our Hopebearer, all hail. To thee our plaint we lift, children of Eve yet in exile. To thee our aspiring, and longing and weeping, lift we from this vale of sorrow. Ah then, Mary, be our intercessor; hither vouchsafe to turn thine eyes compassionate, and look upon us. And Jesus, blessed Offspring of thy womb, O Mother, shew thou to us when earthly exile endeth. O gentle, O loving, O gracious Virgin Mary. 

\begin{responses}

    \V{Pray for us, O Holy Mother of God.}
    \R{That we may be worthy of the promises of Christ.}
    
\end{responses}

Let us pray.

Almighty, everlasting God, who by the co-operation of the Holy Ghost didst prepare the body and soul of the glorious Virgin-Mother Mary to become a dwelling-place meet for thy Son: grant that as we rejoice in her commemoration; so by her fervent intercession we may be delivered from present evils and from everlasting death. Through the same Jesus Christ thy Son our Lord. Who liveth and reigneth with thee, in the unity of the Holy Ghost, ever one God, world without end.

\begin{responses}
    \R{Amen.}
    \V{May help divine \cross be with us all, forever abiding.}
    \R{Amen.}
\end{responses}

\instruct{Next is said the \textsc{Triple Prayer} of the \textsc{Our Father}, \textsc{Hail Mary}, and \textsc{Apostle's Creed}. One can find the \textsc{Our Father} and \textsc{Apostle's Creed} on pp. 66-67 of the BCP. The \textsc{Hail Mary} is reproduced below.}

\instruct{The \textsc{Hail Mary} is composed of the Angelic Salutation to the Blessed Virgin Mary that she would conceive of the Holy Ghost, and where she became the first Christian disciple in \bibleref{Luke 1:28}. It is also composed of her cousin St. Elizabeth's greeting in \bibleref{Luke 1:42}. The final part is a petition for her prayer.}

\begin{prayer}

  Hail Mary, full of grace, the \textsc{Lord} is with thee:
  Blessed art thou among women,
    and blessed is the fruit of thy womb, Jesus.
  Holy Mary, Mother of God, 
    pray for us sinners now, 
    and at the hour of our death. Amen.
    
\end{prayer}

\subsection*{Optional: Prayers After the Office}

\instruct{These are not mandatory, but one can say them for any mistakes that you made during the office.}

\subsubsection*{\textit{Sacrosanctae}}

\blankline
\begin{prayer}

  To \textsc{God} Most Holy, in his Divine Majesty of Trinity in Unity;
  To Jesus Christ, our Lord and God
    made man and crucified for us;
  To blessed Mary Ever-Virgin, from whose
    glorious purity he took flesh;
  And to the entire Company of Saints of God in Heaven;
  Be praise, honor, power, and glory,
    from every creature here on earth.
  And likewise to us sinner may there be
    full remission of all our sins:
  Throughout all ages, world without end. Amen.
  
\end{prayer}

\begin{responses}
    \V{Blessed is the womb of the Virgin Mary, which bore the Son of the everlasting Father.}
    \R{And blessed are the paps which gave suck to Christ our Lord.}
\end{responses}

\instruct{Then is said the \textsc{Our Father} and \textsc{Hail Mary}.}
\newpage
\section{The Psalter}

\end{document}
